\section{Standard di qualità}
Nel corso dell'analisi e della valutazione della qualità dei processi e del software, il gruppo adotta standard internazionali rigorosi per garantire una valutazione oggettiva e conforme alle migliori pratiche globali. L'utilizzo congiunto di questi standard assicura un approccio strutturato, elevata coerenza e affidabilità lungo l'intero ciclo di vita del progetto.
\subsection{Standard ISO/IEC 25010}
Lo standard ISO/IEC 25010 fornisce il modello di riferimento per la definizione e la valutazione della qualità del prodotto software. Esso definisce un quadro completo di caratteristiche e sottocaratteristiche che permettono di misurare in modo accurato attributi chiave quali la funzionalità, l'affidabilità, l'usabilità, l'efficienza, la manutenibilità e la portabilità del sistema sviluppato.
\subsection{Standard ISO/IEC 12207:1995}
Lo standard ISO/IEC 12207:1995 stabilisce l'architettura dei processi del ciclo di vita del software. Tale suddivisione guida l'organizzazione del lavoro del gruppo attraverso la classificazione dei processi in tre macro-categorie fondamentali:
\begin{itemize}
    \item \textbf{Processi primari}: comprendono le attività centrali e dirette alla realizzazione del prodotto, dall'acquisizione dei requisiti fino allo sviluppo e alla manutenzione;
    \item \textbf{Processi di supporto}: affiancano i processi primari garantendo la qualità, la gestione della configurazione, la documentazione, la verifica e la validazione;
    \item \textbf{Processi organizzativi}: definiscono la struttura e il contesto gestionale complessivo all'interno del quale si svolgono le attività di progetto.
\end{itemize}
